%
% teil2.tex -- Beispiel-File für teil2 
%
% (c) 2020 Prof Dr Andreas Müller, Hochschule Rapperswil
%
% !TEX root = ../../buch.tex
% !TEX encoding = UTF-8
%
\section{Der Kreis
\label{elektro:section:DerKreis}}
\kopfrechts{Der Kreis}
Bevor wir im nächsten Kapitel die Transformation von Kreisen betrachten, wollen wir uns mit den Eigenschaften des Kreises beschäftigen.
Der Kreis ist die einfachste und am besten verstandene Figur der Mathematik.
Elektrostatische Feldlinien und auch das Potential können leicht berechnet werden.
Für die elektrostatische Betrachtung dieses Kapitels ist der Kreis optimal, denn sowohl die Formeln für $\vec{E}$ als auch für $\varphi$ hängen nur vom Abstand zum Mittelpunkt $r$ ab.


\subsection{Radialsymmetrische Potentiale
\label{elektro:subsection:RadialsymmetrischePotentiale}}
Im Kreis ist das Potential radialsymmetrisch verteilt.
Der Kreis stellt eine Punktladung mit der Ladungsdichte $\rho$ dar.
%Zur Berechnung des Feldes wird das ... verwendet in seiner Integralform. 
Dank der Punktsymmetrie des Kreises ist das Potentialfeld durch Polarkoordinaten, die den Abstand $r$ und den Winkel $\varphi$ darstellen, vollständig beschrieben.  
Im Abschnitt \ref{elektro:subsection:Einheitskreis-Beispiel} wurde gezeigt wie das Potentialfeld $\varphi$ und das elektrische Feld $\vec{E}$ zusammenhängen.
Es ergibt sich ein Kreislauf von Integralen, in welchem stets von einem Fall auf den anderen geschlossen werden kann und umgekehrt.
Dazu muss jedoch eine dieser Formeln, die für den Kreis gültig ist, aufgestellt worden.
Es braucht also einen Einstieg, durch den die Integralformel des elektrischen Feldes $\vec{E}$ gefunden werden kann. 

\begin{equation}
    \varphi = \int_{-\infty}^{\infty} \frac{\rho}{4\pi\varepsilon_0 r} dV
    \label{elektro:equation:potentialbegin}
\end{equation}
Wie zuvor gezeigt, stellt der Kreis die Ladung dar und wir wollen das Potentialfeld ausserhalb des Kreises beschreiben.
Dazu können wir das Integral von der Kreisoberfläche bis in die Unendlichkeit lösen.
\begin{equation}
    \varphi = \int_{0}^{2\pi} \int_{r}^{\infty} \frac{\rho}{4\pi\varepsilon_0 r} r \, dr \, d\varphi
    = \frac{\rho}{2\varepsilon_0} \int_{r}^{\infty} dr
    = \frac{\rho}{2\varepsilon_0} [r]_{r}^{\infty}
    = -\frac{\rho}{2\varepsilon_0} r
\end{equation}
Die abbildung dieser Formel ist in Abbildung \ref{elektro:figure:phi_feld_kreis} dargestellt.

Verständnis: 
Das Volumenelement wird zu $dV = r \, dr \, d\varphi $.
Für beide Integrale wird die Integrationsgrenze von $0$ bis $2\pi$ für $\varphi$ und von $r$ bis $\infty$ für $r$ gewählt, da das Potentialfeld ausserhalb des Kreises beschrieben werden soll und dies um die gesamten $2\pi$ des Kreises.


\subsection{Kreisspiegelung
\label{elektro:subsection:Kreisspiegelung}}

Damit man sich die Kreisspiegelung vorstellen kann, wird der Kreisrand als Spiegeloberfläche angenommen. 
Somit wird jeder Punkt ausserhalb des Kreises --- auch die Unendlichkeit --- in das Innere des Einheitskreises gespiegelt.
Analog dazu wird alles innerhalb des Einheitskreises nach aussen gespiegelt, wobei der Ursprung auf alle Seiten in die Unendlichkeit gespiegelt wird. 

Diese Transformation wird hier nur gezeigt, weil durch diese Transformation die Wahl der Integralgrenzen im Abschnitt \ref{elektro:subsection:RadialsymmetrischePotentiale} erklärt werden kann.
Es kann vom Mittelpunkt zum Kreisrand oder vom Kreisrand bis in die Unendlichkeit integriert werden und mit einer einfachen Kreisspiegelung im Komplexen vom einen zum anderen Fall übergegangen werden. 
Mathematisch wird die Kreisspiegelung durch die Transformation
\begin{equation}
    z \mapsto \frac{1}{\overline{z}}
\end{equation}
beschrieben.
\textcolor{red}{Kreisspiegelung Fallüberführung rechenbeispiel}


\subsection{Zusammenhang zwischen $\vec{E}$-Feld und Potential
\label{elektro:subsection:ephizusammenhang}}
Anhand des Beispiels des Einheitskreises wird hier gezeigt wie von einem Potentialfeld $\varphi$ auf das elektrische Feld $\vec{E}$ geschlossen werden kann, wie im Abschnitt \ref{elektro:subsection:Einheitskreis-Beispiel} behauptet wird.

Unser Ausgangspunkt ist das elektrische Feld $\vec{E}$, das durch die Umformung zu zylindrischen Koordinaten auf $ \vec{E} = \frac{\rho}{2\pi\varepsilon_0 r} \hat{r} $ vereinfacht werden kann.
Wird dies in die Formel \ref{equation:Potential} eingesetzt, so ergibt sich das Potentialfeld $\varphi$. 
\begin{equation}
    \varphi = - \int_{1}^{k} \frac{\rho}{2\pi\varepsilon_0 r} \hat{r} \cdot d\mathbf{r}
    = - \frac{\rho}{2\pi\varepsilon_0}\vert ln(k) \vert_{1}^{r}
    = - \frac{\rho}{2\pi\varepsilon_0} \operatorname{ln}(r)
    \label{elektro:equation:potential}
\end{equation}
Ist nun das Elektrische Feld bekannt, kann daraus auch das elektrische Potential berechnet werden.
% Kontinuierlich kann auch vom Potentialfeld $\varphi$ auf das elektrische Feld $\mathbf{E}$ geschlossen weerden indem die Formel die sich aus dem soeben ausgeführen Integrals umgestellt wird. 
Der Rückschluss vom Potentialfeld $\varphi$ auf das elektrische Feld $\vec{E}$ ist über die Laplace-Gleichung \ref{elektro:subsection:Laplace} möglich.
Es muss konsistent in zylindrischen Koordinaten gearbeitet werden, damit die Ableitung des Potentials $\varphi$ nach $r$ das elektrische Feld $\vec{E}$ ergibt.

\begin{equation}
    \vec{E} = - \nabla \varphi = - \frac{\partial}{\partial r} \left( - \frac{\rho}{2\pi\varepsilon_0} \operatorname{ln}(r) \right) = \frac{\rho}{2\pi\varepsilon_0 r} \hat{r}
    \label{elektro:equation:efeld}
\end{equation}

Hiermit ist gezeigt, dass die Integrale in einem Kreislauf miteinander verbunden sind und von einem zum anderen geschlossen werden kann.

\subsubsection{Komplexe Darstellung}
Die formeln können nun noch für Komplexe Zahlen tauglich gemacht werden indem der Vektor $\vec{r}$ in Kartesisschen Koordinaten dargestellt wird und danach die Feldanteile Komplex geschrieben werden.
//TODO: nochmals drübergehen bezüglich der Schreibweise und der Herleitung.

\begin{equation}
\begin{aligned}
    \vec{E}
    &=
    \frac{\rho}{2\pi\varepsilon_0 r}\hat{r}
    \\[2mm]
    &=
    \frac{\rho}{2\pi\varepsilon_0 r}
    \left(
    \frac{x}{r}\hat{x}
    +
    \frac{y}{r}\hat{y}
    \right)
    \\[2mm]
    &=
    \frac{\rho}{2\pi\varepsilon_0}
    \left(
    \frac{x}{r^2}\hat{x}
    +
    \frac{y}{r^2}\hat{y}
    \right)
    \\[2mm]
    &=
    \frac{\rho}{2\pi\varepsilon_0}
    \left(
    \frac{x}{x^2+y^2}\hat{x}
    +
    \frac{y}{x^2+y^2}\hat{y}
    \right).
\end{aligned}
\end{equation}
Damit ergeben sich die Feldanteile in $x$- und $y$-Richtung zu:

\begin{equation}
    E_x=
    \frac{\rho}{2\pi\varepsilon_0}
    \frac{x}{x^2+y^2}
    \quad\text{und}\quad
    E_y=
    \frac{\rho}{2\pi\varepsilon_0}
    \frac{y}{x^2+y^2}
\end{equation}

mit 
\begin{equation}
    \mathcal{E}(z)=E_x-iE_y
\end{equation}
wird nun das Feld Komplex geschriebn
\begin{equation}
    \mathcal{E}(z)
    =
    \frac{\rho}{2\pi\varepsilon_0}
    \frac{x-iy}{x^2+y^2}
\end{equation}
Eine Komplexe Zahl ist bekanntlich definiert als $z = x + iy$.

\begin{equation}
\begin{aligned}
    \frac{1}{z}
    &=
    \frac{1}{x+iy}
    \\[2mm]
    &=
    \frac{x-iy}{(x+iy)(x-iy)}
    \\[2mm]
    &=
    \frac{x-iy}{x^2+y^2}.
\end{aligned}
\end{equation}
Wird damit die Formel für $\mathcal{E}(z)$ umgeschrieben, so ergibt sich:
\begin{equation}
    \boxed{
    \mathcal{E}(z)
    =
    \frac{\rho}{2\pi\varepsilon_0}
    \frac{1}{z}
    }
\end{equation}


% \subsection{Cayley-Transformation
% \label{elektro:subsection:Cayley}}
% Dies wäre eine Option aber ist weniger intuitiv, weil alles über die Unendlichkeit geschlossen werden muss.
% Ausserdem, haben die meisten Geometrien eine geschlossene Oberfläche.
